% Part: sets-functions-relations
% Chapter: arithmetization
% Section: reflections

\documentclass[../../../include/open-logic-section]{subfiles}

\begin{document}

\olfileid{sfr}{arith}{ref}
\olsection{काही तात्त्विक चिंतन}

तांत्रिक तपशील एवढेच. पण त्यातून साधले तरी काय?

एक तर, अगदी निर्विवादपणे, त्यातून आपल्याला काही {सुंदर} शुद्ध गणित
मिळाले. शिवाय काही खोल संकल्पनात्मक कामगिरीही झाली. वास्तव संख्या आणि
परिमेय संख्या यांतील निर्णायक फरक संपूर्णता गुणधर्म व्यक्त करतो, हे
ओळखणे ही एक सखोल अंतर्दृष्टी होती. तसेच, डेडेकिंड छेद म्हणून वास्तव
संख्यांची स्पष्ट रचना केल्यामुळे विश्लेषणाच्या विषयाला भक्कम आधार मिळतो.
छेदांपासून नेमके असेच क्षेत्र बनत असल्यामुळे \emph{संपूर्ण क्रमित क्षेत्र}
ही संकल्पना सुसंगत आहे, हे आपल्याला माहीत आहे.

एवढे असूनही, या कामगिरीबद्दलच्या काही शंका उघडपणे मांडल्या पाहिजेत.

पहिली शंका अशी की, वास्तव संख्यांचा परिचित (कदाचित अनंत) दशांश विस्तार
विचारात घेण्यापेक्षा छेदांच्या रूपात त्यांचा विचार करणे \emph{अधिक}
काटेकोर आहे, हे स्पष्ट नाही. दशांश विस्तारांनी वास्तव संख्यांची केलेली
ही ``रचना'' कोशी क्रमिकांद्वारे केलेल्या रचनेस काहीशी मिळतीजुळती आहे;
परंतु प्रत्यक्षात ती सतराव्या शतकाच्या प्रारंभापासूनच गणितज्ञांना मूलतः
माहीत होती (पाहा \olref[cauchy]{sec}). वास्तव संख्यांना संपूर्णता गुणधर्म
आहे याची जाणीव झाल्यामुळे काटेकोरपणात खरी वाढ झाली; वास्तव संख्यांची काही
विशिष्ट संच म्हणून रचना करता येणे हे कदाचित स्वतःपुरते फारसे रोचक नाही.

पूर्णांकांचे किंवा परिमेय संख्यांचे (त्याहून खूप सोपे) अंकगणितीकरण त्या
क्षेत्रांतील काटेकोरपणा वाढवते, हे तर आणखी कमी स्पष्ट आहे. येथे एक साधे
निरीक्षण उपयुक्त ठरेल. नैसर्गिक संख्यांच्या क्रमित जोड्यांचे तुल्यतावर्ग
म्हणून पूर्णांकांची \emph{रचना} केल्यानंतर, मग पूर्णांकांच्या क्रमित
जोड्यांचे तुल्यतावर्ग म्हणून परिमेय संख्यांची रचना केल्यानंतर, आणि मग
परिमेय संख्यांचे संच म्हणून वास्तव संख्यांची रचना केल्यानंतर आपण त्या
रचना लगेचच \emph{विसरून जातो}. विशेषतः, कोणालाही गणिती सिद्धतेत या
रचनांचा \emph{उपयोग} करावासा वाटणार नाही (त्या रचना अपेक्षेप्रमाणेच
वागतात हे सिद्ध करताना मात्र अर्थातच त्यांचा उपयोग करावा लागेल).
नैसर्गिक संख्यांच्या संचांच्या संचांच्या संचांच्या संचांच्या संचांच्या
संचांच्या एखाद्या संचाबद्दल बोलण्यापेक्षा वास्तव संख्येबद्दल थेट बोलणे
खूपच सोपे आहे.
%(कारण वास्तव संख्या परिमेय संख्यांचे संच म्हणून रचल्या; आणि परिमेय संख्या पूर्णांकांच्या क्रमित जोड्यांचे तुल्यतावर्ग, म्हणजे पूर्णांकांच्या संचांच्या संचांचे संच म्हणून रचल्या; इत्यादी.)
%2/3 = [2,3]
%	म्हणजे {<2,3>, ... }
%	{ {{2},{2,3}}, ... }
%
%वास्तव संख्या म्हणजे
%	परिमेय संख्यांचे संच
%
%परिमेय संख्या म्हणजे
%	पूर्णांकांच्या संचांच्या संचांचे संच
%
%पूर्णांक म्हणजे
%	नैसर्गिक संख्यांच्या संचांच्या संचांचे संच
	
या व्याख्या आपल्याला पूर्णांक, परिमेय संख्या किंवा वास्तव संख्या
\emph{नेमक्या काय आहेत}, \emph{सत्तामीमांसक दृष्ट्या सांगायचे तर}, हे
सांगतात, याबद्दल तर
सर्वाधिक शंका आहे. म्हणजे, वास्तव संख्या (उदाहरणार्थ) काही विशिष्ट संच
(संचांच्या संचांचे संच\ldots) \emph{आहेत}, हे संशयास्पद आहे. अशा
दृष्टिकोनासमोरील मुख्य अडथळा हा की ही रचना अनेक वेगवेगळ्या पद्धतींनी करता
आली असती. वास्तव संख्यांच्या बाबतीत खरोखर रोचक रीतीने भिन्न असलेल्या काही
रचना आहेत (पाहा \olref[cauchy]{sec}). पण काही वेगळ्या रचना मिळवण्याची
एक अगदी क्षुल्लक पद्धत अशी: \olref[sfr][rel][ref]{sec} प्रमाणे आपण
क्रमित जोडीची व्याख्या किंचित वेगळी करू शकलो असतो; क्रमित जोडीची ही
पर्यायी संकल्पना वापरली असती, तर आपल्या रचना तितक्याच यशस्वी झाल्या
असत्या, पण शेवटी मिळालेल्या वस्तू वेगळ्या असत्या. त्यामुळे पूर्णांक,
परिमेय संख्या आणि वास्तव संख्या यांच्या अनेक प्रतिस्पर्धी संचसैद्धान्तिक
रचना आहेत. आता पूर्णांक (उदाहरणार्थ) \emph{त्या} संचांऐवजी \emph{हेच}
संच आहेत, असा दावा करणे केवळ मनमानीचे (आणि अडचणीचे) ठरेल.
(\olref[sfr][rel][ref]{sec} प्रमाणे, \citealt{Benacerraf1965} यांनी
प्रसिद्ध केलेल्या युक्तिवादाचे हे एक उदाहरण आहे.)

आणखी एक मुद्दा उपस्थित करायला हवा: आपल्या रचनांमध्ये काहीतरी फारच
\emph{विचित्र} आहे. आपण नैसर्गिक संख्यांपासून सुरुवात केली. मग पूर्णांकांची
रचना केली आणि ``पूर्णांकांचा $0$'', म्हणजे $ \equivrep{0,0}{\Intequiv}$,
याची रचना केली. पण $0 \neq \equivrep{0,0}{\Intequiv}$. खरे तर, आपल्या
रचनांनुसार \emph{एकही} नैसर्गिक संख्या पूर्णांक नाही. पण हे सहजज्ञानाच्या
अगदी विरुद्ध वाटते. प्रत्यक्षात \olref[sfr][set][imp]{sec} मध्ये आपण
फारसा युक्तिवाद न करता $\Nat \subseteq \Rat$ असा दावा केला होता. संख्या
नेमक्या \emph{काय आहेत} हे या रचना सांगत असतील, तर तो दावा उघडपणे असत्य
होता.

मग थोडे मागे हटून पाहिल्यास आपण कुठवर पोहोचलो? अनौपचारिक संचसिद्धान्तात
काम करताना आणि नैसर्गिक संख्या गृहीत धरताना, पूर्णांक, परिमेय संख्या आणि
वास्तव संख्या यांना काही विशिष्ट संच \emph{म्हणून वागवणे} आपल्याला शक्य
होते. त्या अर्थाने, या वस्तूंचे सिद्धान्त संचसिद्धान्तात
\emph{अंतःस्थापित} करता येतात. पण या अंतःस्थापनाचा तात्त्विक आशय अजिबात
सरळसोट नाही.

अर्थात, याने या विषयावरील शेवटचा शब्द सांगितलेला नाही!{} मुद्दा फक्त
एवढाच आहे. वास्तव संख्यांच्या अंकगणितीकरणाला खोल तात्त्विक महत्त्व
\emph{आहे} हे दाखवण्यासाठी आणखी काही \emph{तात्त्विक} युक्तिवाद आवश्यक असेल.

%याव्यतिरिक्त: या संपूर्ण प्रकरणात आपण नैसर्गिक संख्या आपल्याला सहज
%\emph{दिलेल्या} आहेत असे मानले. पण त्यांच्याबाबत आपण काय करू शकतो?
%हा पुढील प्रकरणाचा विषय आहे.

\end{document}
